% Compact, complete insert file: two supplementary tables only.
% Required packages: booktabs,longtable,array,ragged2e,pdflscape,amsmath,natbib
\begin{landscape}
\begingroup
\scriptsize
\setlength{\tabcolsep}{2.4pt}
\renewcommand{\arraystretch}{1.08}
\setlength{\LTleft}{0pt}
\setlength{\LTright}{0pt}

\noindent\textbf{Table S4A. Operational criteria for broad apple lifecycle transitions.}
\begin{longtable}{>{\RaggedRight\arraybackslash}p{5.0cm} >{\RaggedRight\arraybackslash}p{8.0cm} >{\RaggedRight\arraybackslash}p{10.0cm} >{\RaggedRight\arraybackslash}p{3.4cm}}
\toprule
\textbf{Transition} & \textbf{Operational criterion} & \textbf{Quantitative reference} & \textbf{Reference}\\
\midrule
\endfirsthead
\multicolumn{4}{l}{\textit{Table S4A continued}}\\
\toprule
\textbf{Transition} & \textbf{Operational criterion} & \textbf{Quantitative reference} & \textbf{Reference}\\
\midrule
\endhead
\midrule
\multicolumn{4}{r}{\textit{Continued on next page}}\\
\endfoot
\bottomrule
\endlastfoot
NDG00 $\rightarrow$ NDG01 &
Established seedling with sustained juvenile vegetative growth. &
No universal numerical cut-off. &
\citep{ferree2003}\\

NDG01 $\rightarrow$ NDG02 &
Juvenile-to-adult phase transition begins. &
Approximately 77 main-stem nodes. &
\citep{zhang2007,atay2020,liu2025}\\

NDG02 $\rightarrow$ NDG03 &
Reproductive competence and regular flowering are established. &
Approximately 77--122 nodes; adult reference $>122$ nodes. &
\citep{zhang2007,atay2020}\\

NDG03 $\rightarrow$ NDG04 &
Persistent decline in vegetative and reproductive performance. &
No universal numerical cut-off. &
\citep{ferree2003}\\

ODG00 $\rightarrow$ ODG01 &
Planting and orchard establishment are completed. &
Event-defined transition. &
\citep{ferree2003}\\

ODG01 $\rightarrow$ ODG02 &
Regular flowering and managed production are established. &
Approximately sixth year; canopy values are system-specific references. &
\citep{ferree2003,middleton2002}\\

ODG02 $\rightarrow$ ODG03 &
Sustained loss of vigour, productivity or tree health. &
No universal TCSA, yield or canopy cut-off. &
\citep{ferree2003}\\

OAR00 $\rightarrow$ OAR01--04 &
Visible bud swelling after dormancy release. &
Approximately 800--1200 CU as a cultivar-dependent reference. &
\citep{gonzaleznoguer2023,parkes2020,meier1994}\\

OAR01--04 $\rightarrow$ OAR10--13 &
Visible leaf emergence after bud break. &
BBCH09: leaf tips about 5 mm; BBCH10: about 10 mm. &
\citep{meier1994}\\

OAR10--13 $\rightarrow$ OAR30--39 &
Current-season shoot elongation becomes evident. &
No universal LAI or shoot-length boundary. &
\citep{meier1994,wunsche1996}\\

OAR30--39 &
Shoot development progression. &
BBCH35: about 50\% final shoot length; BBCH39: about 90\%. &
\citep{meier1994,wang2020}\\

OAR40--45 $\rightarrow$ OAR50--59 &
First flower opens. &
BBCH61: about 10\% open; BBCH65: about 50\% open. &
\citep{meier1994}\\

OAR50--59 $\rightarrow$ OAR70--79 &
Retained fruitlets enter sustained growth. &
No universal fruit-set percentage. &
\citep{meier1994,serra2016}\\

OAR60--62 &
Flower-bud induction, initiation and differentiation are identified microscopically. &
No universal numerical exit threshold. &
\citep{kofler2019}\\

OAR70--79 $\rightarrow$ OAR80--82 &
Ripening begins as fruit approaches final size. &
BBCH75: about 50\% final fruit size; BBCH81: beginning of ripening. &
\citep{meier1994,malladi2010}\\

OAR80--82 $\rightarrow$ OAR90--93 &
Fruit maturity and leaf senescence are assessed separately. &
No single numerical boundary. &
\citep{meier1994}\\

OAR90--93 $\rightarrow$ OAR00 &
Leaf discoloration and leaf fall progress toward dormancy. &
BBCH95: about 50\% leaves discoloured; BBCH97: all leaves fallen. &
\citep{meier1994,gonzaleznoguer2023}\\

\end{longtable}

\vspace{3mm}
\noindent\textbf{Table S4B. Measurement protocols for quantitative indicators used in the apple lifecycle.}
\begin{longtable}{>{\RaggedRight\arraybackslash}p{5.1cm} >{\RaggedRight\arraybackslash}p{2.3cm} >{\RaggedRight\arraybackslash}p{16.7cm} >{\RaggedRight\arraybackslash}p{2.4cm}}
\toprule
\textbf{Indicator(s)} & \textbf{Stage(s)} & \textbf{Measurement protocol} & \textbf{Reference}\\
\midrule
\endfirsthead
\multicolumn{4}{l}{\textit{Table S4B continued}}\\
\toprule
\textbf{Indicator(s)} & \textbf{Stage(s)} & \textbf{Measurement protocol} & \textbf{Reference}\\
\midrule
\endhead
\midrule
\multicolumn{4}{r}{\textit{Continued on next page}}\\
\endfoot
\bottomrule
\endlastfoot
\multicolumn{4}{l}{\textbf{Environmental, timing and orchard context}}\\
Air temperature, relative humidity, wind speed/direction, rainfall and short-wave radiation &
NDG00; ODG00--03; OAR00--93 &
Air temperature, relative humidity, wind, rainfall and short-wave radiation are recorded with a calibrated orchard weather station following WMO siting and exposure guidance. &
\citep{wmo2024}\\

Photosynthetically active radiation (PAR/PPFD) &
ODG02; OAR10--82 &
PAR/PPFD is measured with a quantum sensor or linear ceptometer positioned above and below the canopy. &
\citep{wunsche1996,middleton2002}\\

Soil temperature &
NDG00; ODG00--03; OAR00--93 &
Soil temperature is recorded with temperature probes installed at fixed root-zone depths. &
\citep{dong2001,wmo2024}\\

Phenological event date, DOY, DAFB and photoperiod &
OAR00--93 &
Phenology is scored from repeated BBCH observations; DOY and DAFB are calculated from the observation date and full-bloom date. &
\citep{meier1994,stoeckli2015}\\

\multicolumn{4}{l}{\textbf{Soil and root-zone}}\\
Composite soil sample &
NDG00; ODG00--03 &
A composite soil sample is prepared by mixing 15--20 subsamples collected from a uniform orchard block at feeder-root depth. &
\citep{wsu_soil_sampling}\\

Soil pH and available nutrient profile (SOC, total/extractable N, Olsen P, exchangeable K and named nutrients) &
NDG00; ODG00--03 &
Soil pH and extractable nutrients are measured from the composite sample using the laboratory extraction specified for each analyte. &
\citep{wsu_soil_sampling}\\

Volumetric soil water content (VWC) &
NDG00; ODG00--03; OAR10--93 &
Volumetric soil water content is measured with calibrated capacitance or TDR probes installed at representative root-zone depths. &
\citep{wsu_soil_moisture}\\

Soil-water potential / tension &
ODG01--03; OAR10--82 &
Soil-water potential is measured with tensiometers or matric-potential sensors installed in good soil contact at root-zone depth. &
\citep{wsu_soil_moisture}\\

Apparent soil electrical conductivity (ECa) &
ODG00--03; OAR10--82 &
Apparent soil electrical conductivity is mapped with a georeferenced electromagnetic or conductivity sensor survey. &
\citep{tsoulias2020}\\

Soil texture, bulk density, porosity, saturated hydraulic conductivity and penetration resistance &
ODG00--03 &
Soil texture, bulk density, porosity, hydraulic conductivity and penetration resistance are measured with standard NRCS field and laboratory procedures. &
\citep{nrcs2024fieldbook}\\

\multicolumn{4}{l}{\textbf{Tree identity, establishment and natural developmental traits}}\\
Scion/rootstock identity and tree age &
ODG00--03 &
Scion, rootstock and planting date are recorded as tree-level metadata; tree age is calculated from the planting date. &
\citep{ferree2003}\\

Graft-union height, row spacing, within-row spacing, planting density and row orientation &
ODG00--03 &
Graft-union height and tree/row spacing are measured with a tape or survey pole; planting density is calculated from spacing and row orientation from field bearing. &
\citep{sun2020,ferree2003}\\

Main-stem node number and flowering-node position &
NDG01--03 &
Main-stem nodes are counted sequentially from a fixed basal node; the node carrying the first flower is recorded with the same numbering convention. &
\citep{liu2025,zhang2007,atay2020}\\

Leaf blade area, length, width, aspect ratio, outline/shape and serration &
NDG01--03 &
Leaf area, length and width are obtained from flattened leaves scanned or photographed with a scale; shape and serration are scored with UPOV descriptors. &
\citep{upov2023apple}\\

Leaf attitude/side-shoot angle, thorn presence/density and adventitious-root presence &
NDG01--03 &
Leaf attitude, side-shoot angle, thorniness and adventitious rooting are scored on fixed plant positions using UPOV apple descriptors. &
\citep{upov2023apple}\\

Leaf epidermal cell area/number, stomatal density, trichome density and leaf thickness &
NDG01--03; OAR10--82 &
Epidermal cell size, stomatal/trichome density and leaf thickness are measured from calibrated microscopic images of fully expanded leaves. &
\citep{jia2024}\\

Leaf mass per area (LMA) / specific leaf weight &
NDG01--03; OAR10--82 &
LMA is calculated from one-sided leaf area and oven-dry mass of the same leaves. &
\citep{migicovsky2018,schechter1992}\\

Phase-related phenolic markers (phloridzin, myricitrin, caffeic acid) &
NDG01--03 &
Phloridzin, myricitrin and caffeic acid are quantified from defined tissues by HPLC using analytical standards. &
\citep{zhang2007}\\

Shoot length, elongation rate and terminal growth cessation &
NDG01--04; OAR30--39; OAR79 &
Shoot length is measured repeatedly from shoot base to tip with a ruler or tape; elongation rate is calculated between dates. &
\citep{wang2020}\\

\multicolumn{4}{l}{\textbf{Root system}}\\
Root length, surface area, mean/maximum diameter and volume &
NDG01; ODG03 &
Washed roots are scanned against a contrasting background and root length, area, diameter and volume are obtained by root-image analysis. &
\citep{liu2019,serrano2023}\\

Primary-root count; fine/medium/coarse root length; root dry mass; root-system width/depth/projected area &
ODG03 &
Excavated roots are counted, separated into diameter classes, scanned for dimensions and oven-dried for dry mass. &
\citep{serrano2023}\\

Rootstock:scion diameter ratio at graft union &
ODG03 &
Rootstock and scion diameters are measured with calipers immediately below and above the graft union and expressed as a diameter ratio. &
\citep{serrano2023}\\

Effective rooting depth and root water potential &
ODG01--03; OAR50--82 &
Effective rooting depth is obtained from depth-resolved root observations; root water potential is measured on sampled roots with the cited pressure-potential method. &
\citep{tsoulias2020}\\

Fine- and coarse-root starch &
OAR01--04 &
Fine and coarse roots are separated by diameter class, dried and ground, and starch is measured after extraction and hydrolysis. &
\citep{mochizuki1956}\\

\multicolumn{4}{l}{\textbf{Tree size, canopy, crop load and productivity}}\\
Trunk circumference and trunk cross-sectional area (TCSA) &
NDG01--04; ODG01--03 &
TCSA is derived from trunk circumference measured with a tape around a smooth trunk surface 30 cm above the graft union: $TCSA=C^2/(4\pi)$. &
\citep{gonzalez2023}\\

Tree height, canopy width, projected area and canopy volume &
ODG01--02; OAR10--82 &
Tree height and orthogonal canopy widths are measured with a tape or survey pole; projected area and canopy volume are calculated from these dimensions. &
\citep{sun2020}\\

Leaf area index (LAI): direct/reconstructed &
ODG02; OAR10--82 &
Direct LAI is total leaf area per tree divided by ground area; leaf area is measured with a leaf-area meter and scaled from sampled leaves/shoots. &
\citep{wunsche1996,middleton2002}\\

Leaf area index (LAI): hemispherical imaging &
ODG02; OAR10--82 &
LAI is estimated from hemispherical canopy images acquired with a CI-110 Plant Canopy Imager and analysed with the CI-110 canopy-analysis software. &
\citep{knerl2018}\\

Leaf area index (LAI): optical canopy analyser &
ODG02; OAR10--82 &
Optical LAI is measured with an LAI-2200C using matched above- and below-canopy readings following the manufacturer's isolated-row protocol. &
\citep{licor_lai2200c}\\

Canopy light interception &
ODG02; OAR10--82 &
Canopy light interception is calculated from above- and below-canopy PAR measured with a linear ceptometer: $1-I_{below}/I_{above}$. &
\citep{middleton2002}\\

Crop load / fruit number per tree or TCSA; blossom-cluster density &
ODG02--03; OAR40--82 &
Crop load is obtained by counting fruit or blossom clusters per tree and, when required, dividing the count by same-tree TCSA. &
\citep{stover2001,serra2016}\\

Annual yield, cumulative yield and yield efficiency &
ODG02--03; OAR82 &
Yield is the total fruit mass harvested per tree; yield efficiency is yield divided by same-tree TCSA. &
\citep{plavcova2022,stover2001}\\

Flowering intensity and return bloom &
ODG02--03; OAR40--59 &
Flowering intensity is obtained from blossom-cluster counts or a fixed bloom scale; return bloom is measured on the same trees the following season. &
\citep{plavcova2022,serra2016}\\

Annual trunk radial increment &
ODG03 &
Annual trunk radial increment is measured from cross-dated increment cores collected at a fixed trunk position. &
\citep{plavcova2022}\\

\multicolumn{4}{l}{\textbf{Bud, dormancy and reserve status}}\\
Bud count and external bud size &
OAR00--04; ODG01--02 &
Buds are counted on a fixed shoot/tree basis and bud dimensions are measured with calipers or scale-referenced images. &
\citep{labuschagne2003,upov2023apple}\\

Bud fresh mass, dry mass and water content &
OAR00--04 &
Bud water content is derived from fresh mass and oven-dry mass measured on the same bud sample. &
\citep{milyaev2024}\\

Chilling accumulation (CH, CU, CP) and chilling requirement &
OAR00--04 &
Chilling is calculated from temperature records with the named model (CH, Utah CU or Dynamic Model CP); chilling requirement uses a stated budbreak endpoint. &
\citep{parkes2020,sapkota2021,milyaev2024}\\

Budbreak percentage under forcing &
OAR00--04 &
Budbreak percentage is measured by forcing dormant shoots under controlled conditions and scoring buds at fixed intervals with a stated BBCH criterion. &
\citep{milyaev2024}\\

Bud starch, soluble sugars, sucrose/glucose/fructose, sorbitol and total carbohydrate reserves &
OAR00--04 &
Bud soluble sugars are extracted and quantified chromatographically or enzymatically; starch is measured from the residual tissue after hydrolysis. &
\citep{sapkota2021,raese1977,sydney_anthrone}\\

\multicolumn{4}{l}{\textbf{Leaf morphology, physiology and nutrients}}\\
Leaf type, fresh/dry mass, area, number, water content, volume and internal gas-space volume &
OAR10--82 &
Leaf fresh mass, area and number are measured on a fixed leaf class; dry mass follows oven drying and water content is derived from fresh and dry mass. &
\citep{schechter1992,lakso1984}\\

SPAD chlorophyll index &
OAR10--82 &
SPAD is measured on recently mature leaves with a SPAD meter using two readings on opposite sides of the midrib. &
\citep{neilsen1995}\\

Leaf chlorophyll a, chlorophyll b, total chlorophyll, a:b ratio and carotenoids &
OAR10--93 &
Leaf chlorophylls and carotenoids are extracted from a fixed leaf mass/area and quantified spectrophotometrically or chromatographically. &
\citep{lauzike2020,tartachnyk2004}\\

Leaf mineral nutrients (N, P, K, Ca, Mg and micronutrients) &
OAR10--93 &
Apple leaf nutrients are measured from recently mature leaves collected using the WSU leaf-tissue sampling protocol and analysed by element-specific laboratory methods. &
\citep{wsu_leaf_tissue}\\

Net photosynthesis, stomatal conductance, transpiration and related gas-exchange metrics &
OAR10--93 &
Net photosynthesis, stomatal conductance and transpiration are measured on defined leaves with an infrared gas analyser under fixed chamber conditions. &
\citep{wunsche1996,lauzike2020}\\

Leaf dark respiration &
OAR90--93 &
Leaf dark respiration is measured as CO$_2$ exchange with an infrared gas analyser after dark acclimation. &
\citep{barden1979}\\

Chlorophyll fluorescence ($F_v/F_m$) &
OAR80--93 &
$F_v/F_m$ is measured on dark-adapted leaves with a PAM fluorometer from minimum and maximum fluorescence. &
\citep{tartachnyk2004}\\

Stem water potential by pressure chamber &
OAR10--82 &
Midday stem water potential is measured with a pressure chamber on a shaded lower-canopy leaf bagged for at least 10 min and measured immediately after excision. &
\citep{ucanr_pressure_chamber}\\

Trunk water potential by embedded microtensiometer &
OAR30--82 &
Trunk water potential is recorded continuously with an embedded microtensiometer installed on a flat trunk surface. &
\citep{blanco2024}\\

Canopy temperature and $T_c-T_a$ &
OAR10--82 &
Canopy temperature is measured with an infrared radiometer or thermal camera and $T_c-T_a$ is calculated from simultaneous air temperature. &
\citep{blanco2024}\\

Continuous trunk diameter and maximum daily trunk shrinkage (MDS) &
OAR30--82 &
Continuous trunk diameter is measured with a linear displacement dendrometer; MDS is the daily maximum minus daily minimum diameter. &
\citep{blanco2024}\\

\multicolumn{4}{l}{\textbf{Flowering, pollination and flower-bud formation}}\\
Flower-cluster count, flowers per cluster, bloom percentage, flowering date and fruit set &
OAR40--70 &
Flower clusters and open flowers are counted on fixed trees/branches; bloom percentage and fruit set are calculated from the same tagged population. &
\citep{stoeckli2015,serra2016,losada2013}\\

Flower / fruitlet cell count and floral microscopy &
OAR40--70 &
Flower or fruitlet cells are counted from fixed, sectioned tissues imaged microscopically along a defined anatomical plane. &
\citep{losada2013,malladi2010}\\

Pollinator visitation rate and flower-standardized visitation &
OAR50--59 &
Pollinator visitation is measured as flower visits recorded during fixed timed observations on a known number of open flowers. &
\citep{gilpin2022}\\

Pollen load on flower visitors &
OAR50--59 &
Pollen load is measured by removing pollen from captured flower visitors and counting grains microscopically. &
\citep{gilpin2022}\\

Pollen-tube progression and controlled pollination &
OAR50--59 &
Pollen-tube progression is measured in controlled-pollinated flowers stained with aniline blue and examined by fluorescence microscopy. &
\citep{losada2013}\\

Pollen viability and in-vitro germination &
OAR40--59 &
Pollen viability is scored by FDA fluorescence and in-vitro germination by the cited sucrose/mineral germination medium. &
\citep{calic2021}\\

Flowering overlap, pollination-to-fertilization interval, stigmatic receptivity and effective pollination period (EPP) &
OAR50--59 &
Flowering overlap is calculated from cultivar bloom windows; receptivity, fertilization interval and EPP are obtained from timed controlled-pollination observations. &
\citep{losada2013,roeder2021}\\

Flower-bud histological stage, meristem diameter/height, appendage number and plastochron &
OAR60--62 &
Flower-bud stage and meristem dimensions are measured from FAA-fixed, paraffin-embedded serial sections examined by light microscopy. &
\citep{kofler2019,foster2003}\\

Bourse-bud phytohormones and defined metabolic analytes &
OAR60--62 &
Bourse-bud hormones are extracted from frozen tissue and quantified by targeted UHPLC-MS/MS using internal standards. &
\citep{milyaev2022}\\

\multicolumn{4}{l}{\textbf{Fruit growth, maturity and quality}}\\
Fruit diameter, fresh mass, growth rate and fruitlet retention/abscission &
OAR70--82 &
Fruit diameter is measured at the equator with digital calipers and fresh mass with a balance; repeated tagged-fruit measurements give growth and retention. &
\citep{malladi2010,greene2013}\\

Fruit cortex cell-layer number, cell diameter, projected cell area, CSSI and CNI &
OAR70--79 &
Fruit cortex cells are measured from calibrated histological images; layer number and cell dimensions are used to calculate CSSI and CNI. &
\citep{malladi2010,ng2013}\\

Fruit biochemical starch concentration &
OAR70--82 &
Fruit starch concentration is measured enzymatically from cortex tissue after removal of soluble sugars. &
\citep{ng2013}\\

Fruit firmness &
OAR80--82 &
Fruit firmness is measured on peeled equatorial sides with an 11.1-mm penetrometer probe following the University of Maryland apple maturity protocol. &
\citep{umd_apple_maturity}\\

Soluble solids concentration (SSC) &
OAR70--82 &
SSC is measured from expressed apple juice with a calibrated refractometer and reported as $^\circ$Brix. &
\citep{umd_apple_maturity}\\

Titratable acidity and fruit juice pH &
OAR70--82 &
Titratable acidity is measured by NaOH titration of a known juice volume; juice pH is measured with a calibrated pH meter. &
\citep{rutkowski2008,umd_apple_maturity}\\

Starch-iodine pattern index and Streif index &
OAR80--82 &
The starch index is measured on an equatorial fruit half stained with iodine--KI and scored against the cultivar chart following the WSU apple harvest protocol. &
\citep{wsu_apple_harvest}\\

Fruit dry mass, dry matter and density &
OAR70--82 &
Fruit dry matter is calculated from fresh and oven-dry mass; density is fruit mass divided by displaced or geometrically measured volume. &
\citep{vieira2022}\\

Background skin colour ($L^*,a^*,b^*$, hue, chroma) &
OAR80--82 &
Background skin colour is measured at fixed fruit positions with a calibrated colorimeter or spectrophotometer in $L^*a^*b^*$ space. &
\citep{felicetti2008}\\

Fruit peel chlorophyll, carotenoids and anthocyanins &
OAR80--82 &
Peel chlorophylls, carotenoids and anthocyanins are extracted from a fixed peel area/mass and quantified spectrophotometrically or by HPLC. &
\citep{felicetti2008}\\

Internal ethylene concentration and ethylene production rate &
OAR80--82 &
Internal ethylene is measured from core-cavity gas by GC-FID; production rate is measured from headspace accumulation of a fruit sealed for a known time. &
\citep{johnston2009}\\

Index of absorbance difference ($I_{AD}$) &
OAR80--82 &
$I_{AD}$ is measured with a DA meter at fixed equatorial positions as $A_{670}-A_{720}$. &
\citep{cocetta2017,umd_apple_maturity}\\

Fruit surface and near-surface tissue temperature &
OAR70--82 &
Fruit surface temperature is obtained from calibrated thermal images; near-surface tissue temperature is measured with a probe at a fixed insertion depth. &
\citep{li2014,wangfst2020}\\

Fruit mineral nutrients (N, P, K, Ca, Mg, B, Cu, Fe, Mn, Zn) &
OAR70--82 &
Fruit mineral nutrients are measured from dried, ground fruit tissue after element-specific digestion using Kjeldahl, AAS or ICP methods. &
\citep{nachtigall2006}\\

Fruit soluble sugars and sorbitol &
OAR70--82 &
Fruit glucose, fructose, sucrose and sorbitol are extracted from defined tissue and quantified by HPLC or enzymatic analysis. &
\citep{filip2016}\\

Postharvest fruit mass loss / water-loss rate &
OAR82; postharvest &
Postharvest water loss is measured by repeatedly weighing the same fruit during storage and expressing mass loss relative to initial mass. &
\citep{johnston2009}\\

\multicolumn{4}{l}{\textbf{IPDM, senescence and end of productive life}}\\
Disease incidence and disease severity &
ODG01--03; OAR10--93 &
Disease incidence is affected units divided by units assessed; severity is scored with the named disease scale following the UC IPM apple protocol. &
\citep{ucipm_apple}\\

Mite eggs/motile mites, aphid colonies, codling-moth trap capture and pest-damaged fruit/russet &
ODG01--03; OAR10--82 &
Mites, aphids, codling moth and fruit damage are monitored with the leaf/shoot/fruit counts and pheromone traps specified in the UC IPM apple protocol. &
\citep{ucipm_apple}\\

Leaf discoloration and leaf fall &
OAR90--93 &
Leaf discoloration and leaf fall are scored repeatedly on the same canopy population using BBCH senescence stages. &
\citep{meier1994}\\

Senescing-leaf photosynthesis, transpiration, stomatal conductance and $F_v/F_m$ &
OAR90--93 &
Gas exchange and $F_v/F_m$ of senescing leaves are measured with the same IRGA and PAM protocols used for mature leaves. &
\citep{tartachnyk2004}\\

Senescing-leaf chlorophyll, mineral nutrients, protein, amino acids, soluble sugars and H$_2$O$_2$ &
OAR90--93 &
Senescing-leaf pigments, nutrients, protein, amino acids, sugars and H$_2$O$_2$ are measured from fixed leaf samples with the cited biochemical assays. &
\citep{spencer1972,nachtigall2006,sapkota2021}\\

Senescence phytohormones (CK, GA, ethylene, SA, ABA, JA) &
OAR90--93 &
Senescence hormones are extracted from frozen leaf tissue and quantified by targeted LC-MS/MS with internal standards. &
\citep{milyaev2022}\\

Stem soluble sugars, starch and total carbohydrates &
OAR70--00 &
Stem soluble sugars are extracted with 80\% ethanol and starch from the residue; both are quantified with the cited anthrone method. &
\citep{sivaci2006,sydney_anthrone}\\

Whole-tree decline severity, survival and mortality &
ODG03 &
Tree decline, survival and mortality are scored repeatedly on the same original tree positions with the cited decline scale. &
\citep{wunsch2024,lawrence2025}\\

Multi-year yield/growth decline &
ODG03 &
Multi-year growth and yield decline are calculated from repeated annual measurements on the same trees. &
\citep{plavcova2022,lawrence2025}\\

Internal trunk necrosis and root-system deterioration &
ODG03 &
Trunk necrosis is measured as necrotic area divided by total cross-sectional wood area; root deterioration is quantified from excavated and scanned roots. &
\citep{serrano2023}\\

\end{longtable}
\endgroup
\end{landscape}
